\section{Einleitung und Historischer Überblick}
Die Entdeckung der Supraleitung durch Heike Kamerlingh Onnes im Jahr 1911 markiert einen fundamentalen Meilenstein der Festkörperphysik. Bei der Untersuchung von Quecksilber stellte er fest, dass der elektrische Widerstand unterhalb einer kritischen Temperatur $T_c$ abrupt auf einen unmessbar kleinen Wert abfällt. Ein Supraleiter ist jedoch nicht nur ein idealer Leiter ($R = 0$). Wie Walther Meißner und Robert Ochsenfeld 1933 zeigten, verhält sich ein Supraleiter im Inneren auch als idealer Diamagnet ($\chi = -1$), der ein von außen angelegtes Magnetfeld vollständig aus seinem Inneren verdrängt (Meissner-Ochsenfeld-Effekt). 

Die mikroskopische Erklärung dieses makroskopischen Quantenphänomens gelang erst 1957 durch John Bardeen, Leon Cooper und Robert Schrieffer (BCS-Theorie). 1962 postulierte Brian D. Josephson schließlich theoretisch, dass supraleitende Ladungsträger (Cooper-Paare) durch eine dünne isolierende Barriere tunneln können, was zur Entdeckung des nach ihm benannten Josephson-Effekte führte. Dieser Versuch gilt der experimentellen Untersuchung dieser Effekte an einer Niob-basierten Josephson-Struktur.
